There is a small logical point to flag first.

From the excerpt of Subproblem \(1\) that you gave, I can rigorously justify parts \( (2)\) and \( (3)\) once an opposite-sign endpoint pair is chosen. But the hint for part \( (1)\) implicitly uses one extra fact that is **not** stated in your excerpt, namely that every translate \(t\in T\) actually occurs as some \(t_s\) (equivalently: every translate contributes at least one of the lonely points). Likewise, the “without loss of generality” in part \( (2)\) implicitly assumes that the chosen pair \((s_a,s_b)\) is an exposed pair of \(\operatorname{conv}(S)\). Under those intended extra facts, the proof is short and complete. I will write it that way. I also use the ambient hypothesis from the usual lonely-point-lemma setup that \(S\subset\mathbb{R}^2\) is in strictly convex position, so every point of \(S\) is a vertex of \(\operatorname{conv}(S)\). ([num.math.uni-goettingen.de](https://num.math.uni-goettingen.de/~hchen/documents/Balancing.pdf))

Let
\[
P:=\operatorname{conv}(S).
\]

## \( (1)\) Opposite signs occur among the lonely points

Assume, for contradiction, that all lonely-point values have the same sign.

If they are all \(+1\), then for every \(s\in S\),
\[
F(\ell_s)=\varepsilon(t_s)=+1.
\]
Under the intended equality-case fact that every translate \(t\in T\) is equal to some \(t_s\), it follows that
\[
\varepsilon(t)=+1\qquad\text{for all }t\in T,
\]
contradicting the hypothesis that \(\varepsilon\) also attains the value \(-1\).

The case where all \(F(\ell_s)=-1\) is identical.

Hence the lonely-point signs are not all the same. Therefore there exist
\[
s_a,s_b\in S
\]
such that
\[
F(\ell_{s_a})=+1,\qquad F(\ell_{s_b})=-1.
\]

## \( (2)\) Boundary-chain decomposition

Now choose the pair \((s_a,s_b)\) as in the statement, and use the intended “without loss of generality” step: choose affine coordinates so that
\[
x(s_a)<x(s)<x(s_b)\qquad\text{for every }s\in S\setminus\{s_a,s_b\}.
\]
Thus \(s_a\) is the unique leftmost vertex of \(P\), and \(s_b\) is the unique rightmost vertex of \(P\).

Remove the two endpoints \(s_a,s_b\) from the boundary \(\partial P\). Since \(P\) is a convex polygon with at least \(3\) vertices, \(\partial P\setminus\{s_a,s_b\}\) has exactly two connected components. Adding the endpoints back gives two polygonal chains from \(s_a\) to \(s_b\); call them \(U\) and \(L\).

### Strict \(x\)-monotonicity
Fix \(c\in\mathbb{R}\). The vertical line \(x=c\) intersects the convex polygon \(P\) in either the empty set or a line segment. Therefore each of the two boundary chains can meet that line in at most one point. Hence each chain is \(x\)-monotone.

Because the \(x\)-direction was chosen generic, no two vertices of \(S\) have the same \(x\)-coordinate. So the monotonicity is actually strict: along either chain, the \(x\)-coordinate strictly increases from \(s_a\) to \(s_b\).

### Convex-arc property
Each chain is a connected subarc of the boundary of the convex polygon \(P\). Therefore each is a convex polygonal arc from \(s_a\) to \(s_b\).

### Counting vertices
The vertex sets of \(U\) and \(L\) cover all of \(S\), and they intersect exactly in the two endpoints:
\[
U\cup L=S,\qquad U\cap L=\{s_a,s_b\}.
\]
Hence
\[
|U|+|L|
=|U\cup L|+|U\cap L|
=|S|+2
=n+2.
\]

This proves part \( (2)\).

## \( (3)\) Each chain contains both signs

By construction, both chains \(U\) and \(L\) contain the endpoints \(s_a\) and \(s_b\). Their signs are
\[
\varepsilon(t_{s_a})=F(\ell_{s_a})=+1,\qquad
\varepsilon(t_{s_b})=F(\ell_{s_b})=-1.
\]
Therefore on each chain, the sign assignment
\[
s\longmapsto \varepsilon(t_s)
\]
takes both values \(+1\) and \(-1\). In particular, it is not constant on either chain.

So part \( (3)\) follows immediately.

---

### Conclusion
Under the two implicit assumptions used by the hint,

1. every \(t\in T\) occurs as some \(t_s\), and  
2. the chosen pair \((s_a,s_b)\) is taken to be an exposed pair of \(P=\operatorname{conv}(S)\),

the three assertions follow exactly as above.

If you want, I can also help with one of two next steps:

- rewriting this into a polished theorem-proof style suitable for your writeup, or
- isolating exactly what additional lemma is needed to make part \( (1)\) completely rigorous from the equality case alone.